% ===========================================================================
% Recomendaciones
% ===========================================================================

\chapter{Recomendaciones}
\label{ch:recomendaciones}

A partir de los resultados obtenidos y las limitaciones identificadas en la presente investigación, se proponen las siguientes recomendaciones para trabajos futuros:

\begin{enumerate}
\item \textbf{Validación experimental.} Ejecutar la simulación Monte Carlo de la implementación mejorada (Exel+AKF) para obtener resultados definitivos, realizando un barrido sistemático de los parámetros del AKF ($\mathbf{Q}_0$, $\mathbf{P}_0$, $N$) para optimizar la relación precisión-\textit{overhead}.

\item \textbf{Implementación en dispositivos reales.} Aplicar la implementación en plataformas inalámbricas reales (Wi-Fi IEEE 802.11, LoRa) para evaluar el comportamiento en condiciones de canal reales con desvanecimiento, interferencia y movilidad.

\item \textbf{Estampado de tiempo por hardware.} Emplear hardware especializado de estampado según IEEE 802.1AS, combinado con el método híbrido Exel+AKF, para alcanzar precisiones en el orden de nanosegundos.

\item \textbf{Extensión a redes multi-salto y BMCA.} Evaluar la implementación en redes con múltiples dispositivos y saltos, incorporando el BMCA. La acumulación de errores y la interacción AKF-BMCA son problemas abiertos.

\item \textbf{Filtros no lineales y no Gaussianos.} Explorar EKF, UKF y Filtros de Partículas para manejar no linealidades y distribuciones de retardo no Gaussianas comunes en entornos NLoS.

\item \textbf{Aprendizaje automático para estimación de asimetría.} Investigar técnicas de ML (redes neuronales ligeras, \textit{reinforcement learning}) para predecir la asimetría a partir de métricas de capa física (RSSI, SNR, PER).

\item \textbf{Integración con 5G-TSN.} Evaluar el método propuesto en el contexto de la convergencia 5G-TSN, donde las asimetrías \textit{uplink/downlink} son inherentes.

\item \textbf{Seguridad en la sincronización.} Investigar la robustez del método frente a ataques maliciosos y desarrollar mecanismos de detección integrados en el AKF mediante análisis de innovaciones.

\item \textbf{Análisis de sensibilidad y optimización.} Realizar un análisis exhaustivo de los parámetros del AKF bajo diferentes condiciones (niveles de asimetría, estabilidad de osciladores, distancias) y aplicar optimización multiobjetivo.

\item \textbf{Documentación y reproducibilidad.} Publicar el código fuente en un repositorio abierto con documentación detallada para facilitar la reproducibilidad y la colaboración.
\end{enumerate}
